\documentclass[11pt]{article}

\usepackage{amsmath, amssymb, amsfonts}
\usepackage[hidelinks]{hyperref}
\usepackage{geometry}
\usepackage{enumitem}
\usepackage{xcolor}
\usepackage{listings}
\geometry{a4paper, margin=1in}

\lstdefinelanguage{Lean}{
  morekeywords={import,open,namespace,section,end,def,theorem,lemma,by,have,show,exact,where,class,instance,let,if,then,else,match,with,intro,fun,forall,exists,noncomputable,private},
  sensitive=true,
  morecomment=[l]{--},
  morecomment=[s]{/-}{-/},
  morestring=[b]"
}

\lstset{
  basicstyle=\ttfamily\small,
  breaklines=true,
  columns=fullflexible,
  frame=single,
  backgroundcolor=\color{gray!5},
  keywordstyle=\color{blue!60!black},
  commentstyle=\color{green!40!black},
  stringstyle=\color{red!50!black}
}

\title{Comments on the Final Paper and Lean Formalisation}
\author{ChatGPT review notes for Dirk Kunert}
\date{\today}

\begin{document}
\maketitle

\section{Overall assessment}

After re-reading both the paper and the Lean code, my assessment is that the
mathematical paper is substantially coherent and that the two period formulas
look correct.  The Lean code is also much stronger than the earlier version,
especially because the unit-aware geometric multiplier layer is now present.

The most important point is not a flaw in the period formula.  Rather, the
paper should describe the boundary of the Lean formalisation carefully.  The
Lean development formalises the residue-combinatorial model with the geometric
multiplier threaded through to the final theorems.  It does not formalise the
full real-geometric strip-to-residue derivation from first principles.

Thus, the final paper should say that the residue-combinatorial core has been
formalised in Lean~4, while the reduction from the real strip geometry to the
residue model is hand-verified in the paper.

\section{Main mathematical result}

The multiset formula appears consistent:
\[
\lambda_{\mathrm{multiset}}=
\begin{cases}
N, & D\nmid N,\\
N/D, & D\mid N,
\end{cases}
\qquad
N=\lfloor\omega\alpha\rfloor+\lfloor\omega\beta\rfloor+1,
\qquad
D=\alpha^2+\beta^2.
\]

The set-valued formula also appears consistent:
\[
\lambda_{\mathrm{set}}=
\begin{cases}
N, & N<D,\\
1, & N\ge D.
\end{cases}
\]

The geometric reduction is correct.  With
\[
r=\alpha x-\beta y,
\]
the strip condition gives
\[
-\beta\omega\le r\le \alpha\omega.
\]
Hence the number of admissible integer values of~\(r\) is
\[
\lfloor\alpha\omega\rfloor+\lfloor\beta\omega\rfloor+1.
\]
The residue formula
\[
\widetilde p\equiv -\alpha\beta^{-1}r\pmod D
\]
is also correct.  Indeed, solving
\[
\beta x+\alpha y=\widetilde p,
\qquad
\alpha x-\beta y=r
\]
gives
\[
x=\frac{\beta\widetilde p+\alpha r}{D},
\qquad
 y=\frac{\alpha\widetilde p-\beta r}{D}.
\]
Integrality of~\(x\) requires
\[
D\mid \beta\widetilde p+\alpha r.
\]
Since \(\beta\) is invertible modulo~\(D\), this gives exactly
\[
\widetilde p\equiv -\alpha\beta^{-1}r\pmod D.
\]

\section{Recommended correction: describe the Lean scope more precisely}

The abstract currently says, in substance, that the algebraic core has been
formalised in Lean~4.  That wording is good.  However, some later passages may
sound as if the complete real-geometric construction has been formalised.
This should be avoided.

I recommend replacing the relevant statement with something like the following:

\begin{quote}
The residue-combinatorial core of both proofs has been formalised in Lean~4.
The formalisation does not formalise the real-geometric strip-to-residue
reduction itself; instead, it starts from the resulting residue model and
verifies the residue distribution, cyclic-interval minimality, degenerate case,
and both multiset and set-valued period theorems, with the geometric multiplier
\(c_r\equiv -\alpha\beta^{-1}r\pmod D\) threaded through to the final
statements.
\end{quote}

This wording is safer and better aligned with the actual Lean code.

\section{Important Lean/paper mismatch: positivity of \texorpdfstring{\(\alpha,\beta\)}{alpha,beta}}

The paper assumes
\[
\alpha,\beta\in\mathbb N=\{1,2,3,\ldots\}.
\]
The Lean code uses
\begin{lstlisting}[language=Lean]
(α β : ℕ) (h_coprime : Nat.Coprime α β)
\end{lstlisting}
but does not globally require
\begin{lstlisting}[language=Lean]
(hα : 0 < α) (hβ : 0 < β)
\end{lstlisting}
in the public-facing theorem statements.

This means that the Lean theorem is formally stated over a slightly wider
domain than the paper.  For example, \(\alpha=0,\beta=1\) satisfies
\texttt{Nat.Coprime α β}, but it is not part of the paper's intended positive
slope setup.

This is not fatal, but I recommend adding explicit positivity assumptions to
the public Lean wrappers, especially
\begin{lstlisting}[language=Lean]
main_theorem_geometric_concrete
set_main_theorem_geometric_concrete
\end{lstlisting}
Even if the proof does not need these assumptions, they make the Lean statement
match the paper.

A possible public theorem signature would include:
\begin{lstlisting}[language=Lean]
(α β : ℕ) (hα : 0 < α) (hβ : 0 < β)
(h_coprime : Nat.Coprime α β)
\end{lstlisting}

Alternatively, the paper can explicitly say that the Lean theorem is stated in
a slightly more general natural-number domain, while the paper itself restricts
to positive \(\alpha,\beta\).

\section{Generic minimality proof}

The proof of the generic case \(D\nmid N\) is now convincing.  The key chain of
reasoning is:

\begin{enumerate}[label=\arabic*.]
    \item \(N\) is a period.
    \item If \(\lambda_0\) is the minimal positive period, then
          \(\lambda_0\mid N\).
    \item A block of \(\lambda_0\) gaps has constant sum \(\sigma\).
    \item Since \(N=k\lambda_0\), one obtains \(k\sigma=D\).
    \item Translation by \(\sigma\) preserves the projected multiset.
    \item Therefore the residue multiplicity function is invariant under
          translation by \(\sigma\).
    \item In the natural multiplier-stepping coordinate, the heavy residue
          classes form a proper cyclic interval.
    \item A proper cyclic interval has trivial translational stabiliser.
    \item Hence no proper divisor of \(N\) can be a period.
\end{enumerate}

This is a solid minimality argument.

One small stylistic recommendation: use \(\lambda_0\) instead of
\(\lambda'\) for the alleged smaller minimal period.  For example:

\begin{quote}
Now suppose, for contradiction, that the minimal positive period
\(\lambda_0\) satisfies \(\lambda_0<N\). Since \(N\) is a period,
Lemma~\ref{lem:minperiod_divides_periods} gives \(\lambda_0\mid N\).
\end{quote}

This avoids the impression that \(\lambda'\) is an arbitrary period rather
than the minimal one.

\section{Set-valued theorem}

The set-valued theorem appears correct.  The dichotomy is exactly the right
one:
\[
N<D \quad\text{versus}\quad N\ge D.
\]
When \(N<D\), every residue class is hit at most once, so the set and multiset
sequences agree.  When \(N\ge D\), every residue class is hit at least once,
so the projected set is all of \(\mathbb Z\) in the \(\widetilde p\)-coordinate,
and the gap sequence is constantly~\(1\).

One wording issue occurs in the introduction.  The text says, in substance,
that when \(N>D\) the set-valued period collapses to~\(1\).  The collapse
already happens at \(N=D\).  I recommend writing:

\begin{quote}
When \(N\ge D\), the set-valued period collapses to~\(1\), because every
integer residue class is then hit.
\end{quote}

The comparison corollary is correct:
\[
\lambda_{\mathrm{set}}\mid \lambda_{\mathrm{multiset}},
\qquad
\lambda_{\mathrm{set}}\le \lambda_{\mathrm{multiset}},
\]
with equality if and only if \(N\le D\).

\section{Lean code: conceptual structure}

The Lean code has three conceptual layers.

\subsection{Untwisted residue model}

The untwisted residue model uses
\begin{lstlisting}[language=Lean]
count_hits D r0 N x
\end{lstlisting}
and the normalised enumeration.  This proves the core residue distribution
and minimality theorem for the case where the admissible residues are simply
\[
r_0,r_0+1,\ldots,r_0+N-1\pmod D.
\]

\subsection{Unit-aware residue model}

The unit-aware model is the important new part.  It uses
\begin{lstlisting}[language=Lean]
count_hits_unit D u r0 N x
\end{lstlisting}
and
\begin{lstlisting}[language=Lean]
difference_sequence_unit α β ω u
\end{lstlisting}
This layer handles multiplication by an arbitrary unit \(u\in(\mathbb Z/D\mathbb Z)^\times\).

\subsection{Geometric specialisation}

The true Lean counterparts of the paper's headline theorems are the geometric
specialisations
\begin{lstlisting}[language=Lean]
main_theorem_geometric_concrete
set_main_theorem_geometric_concrete
\end{lstlisting}
where
\begin{lstlisting}[language=Lean]
u = multiplier α β h_coprime
\end{lstlisting}
This is the layer corresponding to the actual paper residue map
\[
c_r\equiv -\alpha\beta^{-1}r\pmod D.
\]

The older declarations
\begin{lstlisting}[language=Lean]
main_theorem_concrete
set_main_theorem_concrete
\end{lstlisting}
are still useful, but they refer to the normalised residue enumeration, not
directly to the geometric multiplier.  In the paper, the declarations
\texttt{main\_theorem\_geometric\_concrete} and
\texttt{set\_main\_theorem\_geometric\_concrete} should be emphasised as the
final machine-checked statements corresponding to the paper.

\section{Lean code: naming issue}

The typeclass
\begin{lstlisting}[language=Lean]
class GeometricProjection ...
\end{lstlisting}
is somewhat misleadingly named.  It does not encode an actual real-geometric
projection from the strip.  It encodes an abstract residue-combinatorial model
satisfying period and translation-invariance axioms.

A more accurate name would be something like
\begin{lstlisting}[language=Lean]
ResidueProjectionModel
\end{lstlisting}
or
\begin{lstlisting}[language=Lean]
ResidueGapModel
\end{lstlisting}
This is not a mathematical error, but a reviewer reading the code may object
that \texttt{GeometricProjection} is not literally geometric.

If the name is kept, the paper should describe it as an abstract residue-model
interface rather than as a full formalisation of the geometric projection.

\section{Lean correspondence table}

The paper's Lean correspondence table currently includes exact line numbers.
This is fragile, because any small edit to the Lean file will make the line
numbers inaccurate.

I recommend removing the line-number column.  Instead of a table of the form
\[
\text{Paper result} \quad | \quad \text{Lean declaration} \quad | \quad
\text{Line},
\]
use only
\[
\text{Paper result} \quad | \quad \text{Lean declaration}.
\]

For example, replace entries like
\begin{center}
\texttt{coprime\_alpha\_D} --- line 11
\end{center}
by simply
\begin{center}
\texttt{coprime\_alpha\_D}.
\end{center}

This will make the paper more durable.

\section{Bibliography and repository reference}

The repository citation currently uses a 2025 date.  If the paper is compiled
or submitted in 2026 and the repository is active in 2026, consider using a
more explicit GitHub-style citation, for example:

\begin{quote}
Dirk Kunert, \emph{cut-and-project repository}, GitHub repository, accessed
2026.
\end{quote}

Also make sure that the repository is public before submission if the paper
claims that the formalisation is available there.

\section{Recommended changes before submission}

Before submitting the paper, I recommend the following concrete changes:

\begin{enumerate}[label=\arabic*.]
    \item Reword the Lean-scope statements to say ``residue-combinatorial
          core'', not full geometric proof.
    \item Explicitly state that the real-geometric strip-to-residue reduction
          is hand-verified in the paper.
    \item Emphasise
          \texttt{main\_theorem\_geometric\_concrete} and
          \texttt{set\_main\_theorem\_geometric\_concrete} as the final Lean
          counterparts of the two headline theorems.
    \item Add \texttt{0 < α} and \texttt{0 < β} assumptions to the public Lean
          theorem statements, or explain why Lean is stated over a slightly
          wider domain.
    \item Remove exact Lean line numbers from the correspondence table.
    \item Change the introductory phrase ``when \(N>D\)'' to
          ``when \(N\ge D\)'' for the set-valued collapse.
    \item Consider renaming \texttt{GeometricProjection} to something like
          \texttt{ResidueGapModel}, or explain the name carefully.
\end{enumerate}

\section{Final verdict}

The paper is mathematically plausible and internally consistent.  The main
proof strategy is now convincing.  The Lean code substantially supports the
algebraic residue core, including the geometric multiplier.  The remaining
weakness is not the formula itself, but the exact description of the
formalisation boundary: the real-geometric strip-to-residue reduction is still
hand-verified, while the residue distribution and period-minimality arguments
are machine-checked.

With the wording changes above, this version is much closer to submission
quality.

\end{document}
